% Pandoc template for the fed_fraud documents.
% Based on proposal_gen/templates/detailed_proposal_template/detailed_1_classic.tex
% (Classic report): centered title page, table of contents, ruled section
% headings and the Geodesic house footer. Fonts are Latin Modern (lmodern).
% Two layouts: the full report (title page and contents) and, with
% -V brief, a compact one-page-brief header in the same style.
\documentclass[letterpaper,11pt]{article}
\usepackage[T1]{fontenc}
\usepackage[utf8]{inputenc}
\usepackage{lmodern}
\usepackage[letterpaper,top=0.9in,bottom=1.2in,left=1in,right=1in,footskip=0.6in]{geometry}
\usepackage[table]{xcolor}
\usepackage{geodesicbrand}
\usepackage{amsmath,amssymb}
\usepackage{enumitem}
\usepackage{array}
\usepackage{calc}
\usepackage{longtable}
\usepackage{makecell}
\usepackage[most]{tcolorbox}
\usepackage{titlesec}
\usepackage{needspace}
\usepackage{etoolbox}
\usepackage{hyperref}
\hypersetup{colorlinks=true,urlcolor=coolblack,linkcolor=coolblack,filecolor=geoblue,
  pdftitle={Transportation},pdfauthor={\geofirm}}
\urlstyle{same}

\setlist{itemsep=2pt,leftmargin=1.4em}
\setlength{\parindent}{0pt}
\setlength{\parskip}{6pt}
\providecommand{\tightlist}{\setlength{\itemsep}{0pt}\setlength{\parskip}{0pt}}
\setcounter{secnumdepth}{0}

% Section style: house rule under each heading
\titleformat{\section}{\Large\bfseries\color{coolblack}}{\thesection}{0.6em}{}[{\color{coolblack}\titlerule[0.5pt]}]
\titleformat{\subsection}{\large\bfseries\color{geomid}}{\thesubsection}{0.6em}{}
\titleformat{\subsubsection}{\normalsize\bfseries\color{geomid}}{\thesubsubsection}{0.6em}{}
\titlespacing*{\section}{0pt}{10pt}{6pt}

% Tables are written by latex/filters.lua (navy header, banded rows)
\setlength{\LTpre}{6pt}\setlength{\LTpost}{6pt}

\begin{document}

% BRIEF HEADER
\thispagestyle{geodesic}
\begin{minipage}[c]{0.13\textwidth}\includegraphics[width=\linewidth]{icon.png}\end{minipage}\hfill
\begin{minipage}[c]{0.84\textwidth}
  {\textsc{\large \geofirm}}\par\vspace{2pt}
  {\small\color{geogray}Federal fraud expansion\enspace\textperiodcentered\enspace When: \compiletime}
\end{minipage}\par\vspace{6pt}
{\color{coolblack}\rule{\linewidth}{0.3mm}}\par\vspace{2pt}
{\LARGE\bfseries\color{coolblack} Transportation}\par
\vspace{2pt}{\large\color{geomid} One-page brief}\par
\vspace{-2pt}{\color{coolblack}\rule{\linewidth}{0.3mm}}\par\vspace{4pt}
\pagestyle{geodesic}

\textbf{Priority:} 6 of 33 (score 3.7). \textbf{FY2026 outlays:} \$125B
(USAspending). \textbf{As of:} 2026-10-27. Public data only.

\hypertarget{bottom-line}{%
\section{Bottom line}\label{bottom-line}}

DOT's money flows mostly as highway, transit and airport grants to
states, which then hire contractors. Its signature fraud is the
Disadvantaged Business Enterprise (DBE) front: a certified small firm
that passes work and money through to a large contractor. That is the
same pattern as SBA 8(a) and set-aside fronts, so the SBA pipeline
carries over directly. Bid rigging on highway lettings is the second
pattern.

\hypertarget{exposure}{%
\section{Exposure}\label{exposure}}

\begin{itemize}
\tightlist
\item
  Over \$125B a year in outlays, most of it federal-aid construction
  under the 2021 infrastructure law.
\item
  DBE pass-through fraud is a recurring DOT OIG and DOJ case type,
  including criminal convictions and False Claims Act settlements.
\item
  DOT does not report a high-risk improper payment program, so most
  exposure is procurement fraud that payment sampling does not see.
\end{itemize}

\hypertarget{public-data-we-can-use}{%
\section{Public data we can use}\label{public-data-we-can-use}}

\begin{itemize}
\tightlist
\item
  \textbf{USAspending grants and subawards} for FHWA, FTA and FAA.
\item
  \textbf{State Unified Certification Program DBE directories:}
  certified firms, owners, addresses, work codes.
\item
  \textbf{State DOT bid tabulations} (many states publish every bid per
  letting).
\item
  \textbf{SAM.gov, FPDS} for direct DOT contracts.
\item
  \textbf{DOT OIG and DOJ press releases} as labeled cases.
\end{itemize}

\hypertarget{detection-plays}{%
\section{Detection plays}\label{detection-plays}}

\begin{enumerate}
\def\labelenumi{\arabic{enumi}.}
\tightlist
\item
  \textbf{DBE fronts:} certified firms that share addresses, phones or
  officers with their prime contractor, have no equipment or staff
  signals, or work almost only for one prime.
\item
  \textbf{Bid rotation:} bidders who take turns winning, submit
  complementary high bids, or share ownership across ``competing''
  firms.
\item
  \textbf{Subaward anomalies:} DBE participation claimed on bids that
  never shows up in subaward reporting.
\end{enumerate}

\hypertarget{first-prototype}{%
\section{First prototype}\label{first-prototype}}

A DBE integrity screen for two or three states with good public
directories and bid tabs. Reuse the SBA entity-resolution code; rank
firms with reasons; validate against charged DBE cases.

\hypertarget{limits}{%
\section{Limits}\label{limits}}

Prime-to-subcontractor payments are reported to states, not publicly,
and data quality varies by state. Results point to firms worth review,
not proof.

\hypertarget{who-to-take-it-to}{%
\section{Who to take it to}\label{who-to-take-it-to}}

DOT OIG, FHWA Office of Civil Rights, state DOT DBE offices, and DOJ's
Procurement Collusion Strike Force.

\hypertarget{sources}{%
\section{Sources}\label{sources}}

\begin{itemize}
\tightlist
\item
  \href{https://www.justice.gov/usao-edpa/pr/downingtown-man-convicted-exploiting-usdot-disadvantaged-business-program}{DOJ:
  Pennsylvania man convicted of exploiting the DOT DBE program}
\item
  \href{https://www.justice.gov/opa/pr/ohio-construction-firm-agrees-pay-500000-resolve-false-claims-act-allegations}{DOJ:
  Ohio construction firm pays \$500K over DBE False Claims Act
  allegations}
\item
  \href{https://www.troutman.com/insights/suppliers-beware-us-government-continues-prosecution-of-dbe-fraud-cases-involving-supplies-passed-through-dbes/}{Troutman:
  continued prosecution of DBE pass-through cases}
\item
  \href{https://www.usaspending.gov/agency}{USAspending agency profiles
  (FY2026 outlays)}
\end{itemize}

\end{document}
